\documentclass[12pt, letterpaper, parskip = full]{article}
\usepackage[mmddyyyy]{datetime}

\usepackage[utf8]{inputenc}
\usepackage{amsmath}
\usepackage{mathpazo}
\usepackage{eulervm}
\usepackage[skip=10pt plus1pt]{parskip}

\usepackage[margin=1in]{geometry}
\linespread{1.20}
\usepackage{enumerate}
\usepackage{enumitem}
\usepackage{listings}
\usepackage[T1]{fontenc}
\usepackage{amsfonts}
\usepackage{amssymb}
\usepackage{amsthm}
\usepackage{multirow,array}
\usepackage{graphicx}
\usepackage{physics}
\usepackage{bm}
\usepackage{tikz}
\usepackage{breqn}
\usepackage[longnamesfirst]{natbib}
\usepackage{comment}
\usepackage{setspace}
\usepackage{hyperref}
\usepackage{accents}
\usepackage{graphicx}

\graphicspath{{Images/}}
\setlength{\parindent}{0pt}

\newtheorem{theorem}{Theorem}
\newtheorem{corollary}{Corollary}
\newtheorem{lemma}{Lemma}
\newtheorem{definition}{Definition}
\newtheorem{axiom}{Axiom}
\newcommand{\e}{\mathrm{e}}
\newcommand{\Var}{\mathrm{Var}}
\newcommand{\Ex}[1]{\mathbb{E}\left[#1\right]}
\renewcommand{\Pr}[1]{\mathbb{P}(#1)}
\newcommand{\plim}{\mathrm{plim}}
\newcommand{\pc}[1]{\frac{\dot{#1}(t)}{#1(t)}}
\renewcommand\qedsymbol{$\blacksquare$}
\newcommand{\ubar}[1]{\underaccent{\bar}{#1}}

\title{Worksheet 7 \\ Actuarial Corporate Finance}
\author{Alejandro Flores-Brown}

\begin{document}

    \maketitle

    \section{}

    We are given a Sharpe ratio of \(0.4\) for the market, as well as a Sharpe ratio of \(0.33\) for the QRS stock, a correlation of \(0.8\) of the stock with the market, and a volatility of \(0.3\). We want to find the alpha for the stock. The equation for \(\alpha\) is given by
    
    \[ \alpha_{QRS} = \Ex{R_{QRS}} - r_{QRS} \]

    where the return \(r_{QRS}\) is given by 

    \[ r_{QRS} = r_f + \beta \left( \Ex{R_M} - r_f \right) \]

    The Sharpe ratio is given by \(\frac{\Ex{R_{QRS}} - r_f}{\sigma_{QRS}}\) so we can write

    \[ 0.33 = \frac{\Ex{R_{QRS}} - r_f}{\sigma_{QRS}}\]

    We can find the expected return by using the correlation to find the \(\beta\) for the stock, which will then yield the expected return via the CAPM. We have \(\beta_{QRS} = \rho \frac{\sigma_{QRS}}{\sigma_{M}} = 0.8 * \frac{0.3}{\sigma_M} \), but this still needs a volatility for the market. But by the market's Sharpe ratio we know 

    \[ 0.4 = \frac{\Ex{R_M} - r_f}{\sigma_M}  \]

    Substituting everything, we can write

    \[  \alpha_{QRS} = \Ex{R_{QRS}} - r_f - \beta \left( \Ex{R_M} - r_f \right) = \Ex{R_{QRS}} - r_f - \rho \frac{\sigma_{QRS}}{\sigma_M} (\Ex{R_M} - r_f)  \]

    Looking at this, we can see that the right-hand quantity contains the Sharpe ratio of the market! The first term of the right-hand side can also be transformed into a Sharpe ratio if we divide through by \(\sigma_{QRS}\):

    \[ \frac{\alpha}{\sigma_{QRS}} = \frac{\Ex{R_{QRS}} - r_f}{\sigma_{QRS}} - \rho SR_M = SR_{QRS} - \rho SR_M \]

    Now we can substitute numbers:

    \[ \frac{\alpha}{0.3} = 0.33 - 0.8 \times 0.4 \implies \boxed{\alpha = 0.003}   \]

    \section{}

    We are given a stock price of \(60\) and future earnings of \(10\%\). No other fundamentals have changed. What happens to the stock price? We can set the stock price as 

    \[ S_0 = PV(Div_{t})  \]

    where the right-hand side is the present value of future dividends. Assuming the dividend is a perpetuity (i.e. dividend payments happen forever), dividends across time are constant,  and no probability of the company failing, we have \(PV(\sum(Div_t)) = \frac{Div}{r}\), which gives

    \[ S_1 = \frac{Div}{r} = S_0 \]

    If we then assume dividends go up by \(1.1\) in conjunction with the future earnings increasing, we then have 
     
    \[ S_1 = \frac{1.1 Div}{r} = 1.1 S_0  \]

    so the stock price goes up \(10\%\) to \(\boxed{S_1 = 66}\).

    \section{}

    Suppose \(ABC\) and \(XYZ\) are both expected to pay \(2.2\) million dividends into perpetuity (per year). \(ABC\)'s cost of capital is \(12\%\); \(XYZ\)'s is \(16\%\). We recall that we could set the return equal to the cost of capital.

    \subsection{}

    Which has the higher market value? We write 
    \(\frac{D}{r_i} = V_i\)
    so it immediately follows that since \(\frac{D}{r_1} = \frac{D_{ABC}}{0.12}\) is higher, it must have higher market value \(V_1\).

    \subsection{}

    As said before, firm \(XYZ\) should have higher expected return since it is equal to the cost of capital.

    \subsection{}

    By definition we know 

    \[ \alpha_s = \Ex{R_s} - r_s \]

    which means that XYZ has the higher \(\alpha\) since \(r_{XYZ} = 16\% - 15\% = 1\%\) while \(\alpha_{ABC} = -3\%\).
    
    \subsection{}

    The logic of the previous problems immediately implies that a higher \(\alpha\) implies a lower market value, since market value is inversely related to expected return.



\end{document}